\textbf{Proof.}
Let \(N=2n\), index exposure coordinates by \(j\), and put \(Z_j=D_j-\pi_j\).  Known exposure and the signed convention in \(Y\) give
\[
\mathbb E_\nu\widehat Y_j^{\mathrm{obs}}=Y_j,
\qquad
\{\Pi^{-1}(\widehat Y^{\mathrm{obs}}-Y)\}_j
=\frac{Y_jZ_j}{\pi_j^2}.
\tag{1}
\]
If \(S_B=0\), its spectral pseudoinverse is zero.  Both coefficients vanish, and substitution into the estimator gives exactly the uncontrolled Horvitz--Thompson estimator.  This proves the rank-zero assertions.

Fix positive rank and set
\[
G_B=\widehat\gamma^{\mathrm{HT}}(B)-\gamma_B^{*,+}(Y).
\]
Equation (1) yields
\[
G_B=\sum_{j=1}^{N}u_{B,j}Z_j,
\qquad
u_{B,j}=S_B^+B^{\top}\Omega e_j\frac{Y_j}{\pi_j^2}.
\tag{2}
\]
The stratum margin implies
\[
\lVert S_B^+\rVert_{\mathrm{op}}\le(n\lambda_r)^{-1}.
\tag{3}
\]
Each covariance column has at most \(s_n\) nonzero entries, covariance entries have absolute value at most one, and basis coordinates are bounded by \(\overline a\).  Therefore
\[
\sup_j\lVert B^{\top}\Omega e_j\rVert_2
\le\sqrt2\,\overline a s_n,
\qquad
\sup_j\lVert u_{B,j}\rVert_2\le U_{n,r}.
\tag{4}
\]
The centered basis statistic is
\[
X_B=n^{-1}B^{\top}(D-\mathbb E_\nu D)
=\sum_{j=1}^{N}x_{B,j}Z_j,
\qquad
x_{B,j}=n^{-1}B^{\top}e_j,
\tag{5}
\]
so \(\sup_j\lVert x_{B,j}\rVert_2\le X_n\).  Applying the fourth-moment envelope coordinatewise gives
\[
\mathbb E_\nu\lVert G_B\rVert_2^4
\le4\mathfrak M_{4,n}U_{n,r}^4,
\qquad
\mathbb E_\nu\lVert X_B\rVert_2^4
\le4\mathfrak M_{4,n}X_n^4.
\tag{6}
\]
The feasible and oracle estimators satisfy
\[
\widehat\tau_B^{\mathrm{CV}}
=\widehat\tau_B^{\mathrm{or},+}+\Delta_B,
\qquad
\Delta_B=-G_B^{\top}X_B.
\tag{7}
\]
Cauchy--Schwarz and (6) imply
\[
\mathbb E_\nu\Delta_B^2
\le4\mathfrak M_{4,n}U_{n,r}^2X_n^2.
\tag{8}
\]

Put \(T_B=\widehat\tau_B^{\mathrm{or},+}-\tau\).  The Horvitz--Thompson term is unbiased by (1), while the basis statistic is centered, so \(\mathbb E_\nu T_B=0\).  Moreover
\[
T_B=\sum_{j=1}^{N}t_{B,j}Z_j,
\qquad
t_{B,j}=n^{-1}\left\{\frac{Y_j}{\pi_j}
-B_{j\cdot}\gamma_B^{*,+}(Y)\right\}.
\tag{9}
\]
Writing \(v=\Omega^{1/2}\Pi^{-1}Y\) and \(Z_B^\circ=\Omega^{1/2}B\), the projector inequality gives
\[
\gamma_B^{*,+}(Y)^{\top}S_B\gamma_B^{*,+}(Y)
\le Y^{\top}\Pi^{-1}\Omega\Pi^{-1}Y.
\tag{10}
\]
There are at most \(2ns_n\) nonzero entries in the covariance sum, so
\[
Y^{\top}\Pi^{-1}\Omega\Pi^{-1}Y
\le2ns_n\pi_{\min,n}^{-2}.
\tag{11}
\]
Equations (10)--(11) and the positive-eigenvalue margin show
\[
\lVert\gamma_B^{*,+}(Y)\rVert_2\le\Gamma_{n,r},
\qquad
\sup_j|t_{B,j}|\le T_{n,r}.
\tag{12}
\]
Expanding the covariance in the two coefficient coordinates and using the third-moment envelope gives
\[
|\operatorname{Cov}_\nu(T_B,\Delta_B\mid Y)|
\le2\mathfrak M_{3,n}T_{n,r}U_{n,r}X_n.
\tag{13}
\]
Since
\[
\operatorname{Var}(T_B+\Delta_B)-\operatorname{Var}(T_B)
=\operatorname{Var}(\Delta_B)+2\operatorname{Cov}(T_B,\Delta_B),
\]
equations (8) and (13) prove the finite bound \(R_{n,r}^{\mathrm{ad}}\).

For the oracle identity, put \(q_B=B^{\top}\Omega\Pi^{-1}Y\).  Since
\[
q_B\in\operatorname{range}\{(\Omega^{1/2}B)^{\top}\}
=\operatorname{range}(S_B),
\]
one has \(S_BS_B^+q_B=q_B\).  Completing the quadratic variance expression gives
\[
\operatorname{Var}_\nu(\widehat\tau_B^{\mathrm{or},+}\mid Y)
=n^{-2}Y^{\top}\left\{\Pi^{-1}\Omega\Pi^{-1}
-\Pi^{-1}\Omega B S_B^+B^{\top}\Omega\Pi^{-1}\right\}Y.
\tag{14}
\]
Under the i.i.d. pair law, \(\mathbb E(YY^{\top})=Q(c_H)\); taking the trace expectation in (14) proves \(\ell_B^+(c_H)\).  Pilot independence permits conditioning on a selected class before applying the same calculation and then integrating.

For the dependency rates, a centered mixed moment vanishes whenever one distinct coordinate is separated in \(\mathcal G_n\) from all remaining coordinates.  Counting ordered contributing triples and quadruples and using \(|Z_j|\le1\) gives
\[
\mathfrak M_{3,n}\le8(2n)d_n^2,
\qquad
\mathfrak M_{4,n}\le64(2n)^2d_n^2.
\tag{15}
\]
Nonadjacent coordinates are independent and hence uncorrelated, so \(s_n\le d_n\).  Substitution of the displayed exposure and dependency rates into \(R_{n,r}^{\mathrm{ad}}\) proves the two-term rate.  Because \((7/2)\alpha+3\beta\le4\alpha+4\beta\), the first-order condition implies \(nR_{n,r}^{\mathrm{ad}}\to0\) uniformly across the two positively separated ranks.  At rank two the stratum is the ordinary \(\kappa\)-regular class and the pseudoinverse is the inverse, so the existing rank-two bridge applies.

Finally, every \(S(b_1,b_0)\) is positive semidefinite.  Such a \(2\times2\) matrix has rank two exactly when its determinant is positive, rank one exactly when its determinant is zero and trace positive, and rank zero exactly when its trace is zero.  Taking existential maxima proves the rank-two and rank-one certificates.  If both diagonal blocks of \(\Omega\) vanish, positive semidefiniteness forces the cross block to vanish; conversely a nonzero diagonal block has a vector with positive quadratic form.  This proves the rank-zero certificate and completes the proof. \(\square\)